\documentclass{beamer}

\usepackage[utf8]{inputenc}
\usepackage[T5]{fontenc}
\usepackage[vietnamese]{babel}

\usepackage{amsmath,amssymb}
\usepackage{tikz}
\usepackage{listings}
\usepackage{array}
\usepackage{colortbl}
\usepackage{booktabs}
\usepackage{xcolor}
\usepackage{tcolorbox}

\usetheme{Madrid}
\usecolortheme{dolphin}

% --- màu ---
\definecolor{myblue}{RGB}{30,144,255}
\definecolor{mygreen}{RGB}{60,179,113}
\definecolor{myred}{RGB}{220,20,60}
\definecolor{mypurple}{RGB}{138,43,226}
\definecolor{myorange}{RGB}{255,140,0}
\definecolor{mycyan}{RGB}{0,206,209}
\definecolor{codebg}{RGB}{248,250,252}
\definecolor{codeid}{RGB}{255,140,0}

\setbeamercolor{structure}{fg=myblue}
\setbeamercolor{frametitle}{fg=white,bg=myblue}
\setbeamerfont{frametitle}{series=\bfseries}

\tcbset{
colback=blue!5,
colframe=myblue,
arc=3mm,
boxrule=0.8pt
}

% macros
\newcommand{\hlblue}[1]{\textcolor{myblue}{#1}}
\newcommand{\hlgreen}[1]{\textcolor{mygreen}{#1}}
\newcommand{\hlred}[1]{\textcolor{myred}{#1}}
\newcommand{\hlpurple}[1]{\textcolor{mypurple}{#1}}
\newcommand{\hlcyan}[1]{\textcolor{mycyan}{#1}}
\newcommand{\hlwhite}[1]{\textcolor{white}{#1}}

% code
\lstset{
basicstyle=\ttfamily\small,
backgroundcolor=\color{codebg},
keywordstyle=\color{myblue}\bfseries,
commentstyle=\color{mygreen}\itshape,
stringstyle=\color{mypurple},
identifierstyle=\color{codeid},
numbers=left,
numberstyle=\tiny\color{gray},
breaklines=true,
showstringspaces=false,
frame=single,
rulecolor=\color{myblue}
}

\title{\hlred{Bellman-Ford Algorithm}}
\author{Group 3}
\date{\hlcyan{03/2026}}

\begin{document}

\frame{\titlepage}

% ===================== KHÁI NIỆM =====================

\begin{frame}{\hlwhite{Định nghĩa}}
\begin{tcolorbox}[title=\textbf{\hlwhite{Khái niệm}}]
\hlblue{Bellman-Ford} là thuật toán tìm:

\[
\text{Shortest Path từ } s \rightarrow v
\]

trong đồ thị có thể chứa \hlred{trọng số âm}
\end{tcolorbox}

\begin{itemize}
\item Áp dụng cho \hlgreen{single-source}
\item Xử lý được cạnh âm
\end{itemize}
\end{frame}

% ===================== Ý TƯỞNG =====================

\begin{frame}{\hlwhite{Ý tưởng}}
\begin{itemize}
\item Lặp qua tất cả các cạnh
\item Thực hiện \hlblue{relaxation}
\end{itemize}

\begin{tcolorbox}
\centering
\hlpurple{Lặp $n-1$ lần}
\end{tcolorbox}

\begin{itemize}
\item Vì đường đi dài nhất có tối đa $n-1$ cạnh
\end{itemize}
\end{frame}

% ===================== RELAX =====================

\begin{frame}{\hlwhite{Relaxation}}
\begin{tcolorbox}
\centering
\[
dist[v] = \min(dist[v], dist[u] + w(u,v))
\]
\end{tcolorbox}

\begin{itemize}
\item Cập nhật khoảng cách tốt hơn
\item Áp dụng cho mọi cạnh
\end{itemize}
\end{frame}

% ===================== NEGATIVE CYCLE =====================

\begin{frame}{\hlwhite{Negative Cycle}}
\begin{itemize}
\item Sau $n-1$ lần:
\begin{itemize}
\item Nếu còn relax được → có chu trình âm
\end{itemize}
\end{itemize}

\begin{tcolorbox}
\centering
\hlred{Negative Cycle → không tồn tại đường đi ngắn nhất}
\end{tcolorbox}
\end{frame}

% ===================== THUẬT TOÁN =====================

\begin{frame}{\hlwhite{Thuật toán}}
\begin{itemize}
\item Khởi tạo:
\begin{itemize}
\item $dist[s] = 0$
\item còn lại = $\infty$
\end{itemize}

\item Lặp $n-1$ lần:
\begin{itemize}
\item Duyệt tất cả cạnh
\end{itemize}

\item Lần thứ $n$:
\begin{itemize}
\item Kiểm tra chu trình âm
\end{itemize}
\end{itemize}
\end{frame}

% ===================== CODE =====================

\begin{frame}[fragile]{\hlwhite{Code (Core)}}
\begin{lstlisting}[language=C++]
for(int i = 1; i <= n-1; i++){
    for(auto e : edges){
        int u = e.u, v = e.v, w = e.w;
        if(dist[u] + w < dist[v]){
            dist[v] = dist[u] + w;
        }
    }
}
\end{lstlisting}
\end{frame}

\begin{frame}[fragile]{\hlwhite{Detect Negative Cycle}}
\begin{lstlisting}[language=C++]
for(auto e : edges){
    int u = e.u, v = e.v, w = e.w;
    if(dist[u] + w < dist[v]){
        // có negative cycle
    }
}
\end{lstlisting}
\end{frame}

% ===================== ĐỘ PHỨC TẠP =====================

\begin{frame}{\hlwhite{Độ phức tạp}}
\begin{itemize}
\item Mỗi lần: duyệt $m$ cạnh
\item Lặp $n-1$ lần
\end{itemize}

\begin{tcolorbox}
\centering
\hlred{$O(n \cdot m)$}
\end{tcolorbox}
\end{frame}

% ===================== SO SÁNH =====================

\begin{frame}{\hlwhite{So sánh}}
\begin{itemize}
\item Dijkstra:
\begin{itemize}
\item Nhanh hơn
\item Không xử lý được cạnh âm
\end{itemize}

\item Bellman-Ford:
\begin{itemize}
\item Chậm hơn
\item \hlgreen{Xử lý được cạnh âm + phát hiện cycle}
\end{itemize}
\end{itemize}
\end{frame}

% ===================== KẾT =====================

\begin{frame}{\hlwhite{Kết luận}}
\begin{itemize}
\item Dùng khi có cạnh âm
\item Quan trọng để detect negative cycle
\item Là nền tảng cho SPFA
\end{itemize}
\end{frame}

\end{document}